\documentclass[12pt,a4paper]{article}

\usepackage[margin=1in]{geometry}
\usepackage{amssymb}
\usepackage{booktabs}
\usepackage{hyperref}
\usepackage{enumitem}
\usepackage{titlesec}
\usepackage{fancyhdr}
\usepackage{xcolor}
\usepackage{array}
\usepackage{tabularx}
\usepackage{multirow}
\usepackage{amsmath}

% ── Colours ───────────────────────────────────────────────
\definecolor{headblue}{HTML}{1A5276}

% ── Section styling ───────────────────────────────────────
\titleformat{\section}
  {\Large\bfseries\color{headblue}}
  {\thesection.}{0.5em}{}[\vspace{2pt}\titlerule]
\titleformat{\subsection}
  {\large\bfseries\color{headblue!80}}
  {\thesubsection}{0.5em}{}

% ── Header / Footer ──────────────────────────────────────
\pagestyle{fancy}
\setlength{\headheight}{14.5pt}
\addtolength{\topmargin}{-2.5pt}
\fancyhf{}
\lhead{\small\textit{Cost Estimation --- CA Firm System}}
\rhead{\small\textit{Group E-17}}
\cfoot{\thepage}

\begin{document}

% ═══════════════════════════════════════════════════════════
% TITLE PAGE
% ═══════════════════════════════════════════════════════════
\begin{titlepage}
\centering
\vspace*{2cm}
{\Huge\bfseries\color{headblue} Software Cost Estimation\\[12pt]}
{\LARGE Function Point Analysis \& Use Case Points\\[8pt]}
{\large CA Firm Attendance and Article Management System\\[20pt]}
\rule{0.6\textwidth}{1pt}\\[20pt]

{\large\textbf{Assignment No. 8}}\\[30pt]

\begin{tabular}{ll}
\textbf{Group:} & E-17 \\[4pt]
\textbf{Members:} & Nivesh Jain \\
 & Harsh Nayak \\
 & Rachit Ingole \\
 & Hasan Rupawalla \\
\end{tabular}

\vfill
{\small Department of Computer Engineering\\
\today}
\end{titlepage}

% ═══════════════════════════════════════════════════════════
\section{Problem Statement}
% ═══════════════════════════════════════════════════════════
Perform software cost estimation exercises using \textbf{Function Point Analysis (FPA)} and \textbf{Use Case Points (UCP)} for the CA Firm Attendance and Article Management System. Estimate the software size, effort, and approximate development cost based on these two industry-standard techniques.

% ═══════════════════════════════════════════════════════════
\section{Introduction}
% ═══════════════════════════════════════════════════════════
Software cost estimation is a fundamental activity in project management that predicts the resources, time, and budget required to develop a software system. Accurate estimation helps prevent budget overruns, enables proper resource allocation, and sets realistic delivery expectations.

Two widely used estimation techniques are:
\begin{enumerate}[noitemsep, label=\roman*)]
    \item \textbf{Function Point Analysis (FPA)} --- measures software size based on the functional requirements delivered to the user, independent of the programming language or technology used.
    \item \textbf{Use Case Points (UCP)} --- estimates effort based on the complexity of use cases and the environmental and technical factors of the project.
\end{enumerate}

This assignment applies both techniques to the \textbf{CA Firm Attendance and Article Management System}, a web-based platform built using React.js, Node.js/Express, and MySQL that automates attendance tracking (biometric face recognition + GPS verification), task assignment, progress monitoring, leave management, and reporting for a Chartered Accountant firm (SPCM).

% ═══════════════════════════════════════════════════════════
% PART A: FUNCTION POINT ANALYSIS
% ═══════════════════════════════════════════════════════════
\section{Function Point Analysis (FPA)}

\subsection{Overview}
Function Point Analysis, introduced by Allan Albrecht at IBM in 1979, quantifies software functionality from the user's perspective. It classifies system interactions into five categories:
\begin{itemize}[noitemsep]
    \item \textbf{External Inputs (EI)} --- data entering the system from outside.
    \item \textbf{External Outputs (EO)} --- data or reports leaving the system.
    \item \textbf{External Inquiries (EQ)} --- input/output combinations for data retrieval.
    \item \textbf{Internal Logical Files (ILF)} --- data maintained within the system.
    \item \textbf{External Interface Files (EIF)} --- data referenced from external systems.
\end{itemize}

\subsection{Step 1: Identify and Classify Function Points}

\subsubsection*{External Inputs (EI)}
\begin{center}
\begin{tabular}{clc}
\toprule
\textbf{\#} & \textbf{Function} & \textbf{Complexity} \\
\midrule
1 & Admin Login & Low \\
2 & Article Trainee Login & Low \\
3 & Add Article Trainee (Admin) & Average \\
4 & Edit Article Trainee Details (Admin) & Average \\
5 & Mark Attendance (Face + GPS) & High \\
6 & Create Task (Admin) & Average \\
7 & Assign Task to Article(s) & Average \\
8 & Update Task Progress (Article) & Average \\
9 & Submit Leave Request (Article) & Low \\
10 & Approve/Reject Leave (Admin) & Low \\
11 & Configure System Settings (Admin) & Low \\
\bottomrule
\end{tabular}
\end{center}

\subsubsection*{External Outputs (EO)}
\begin{center}
\begin{tabular}{clc}
\toprule
\textbf{\#} & \textbf{Function} & \textbf{Complexity} \\
\midrule
1 & Attendance Report (Daily/Monthly) & High \\
2 & Performance Report & High \\
3 & Task Completion Statistics & Average \\
4 & Leave Summary Report & Low \\
5 & Analytics Dashboard (Chart.js) & High \\
6 & Email Notifications (Nodemailer) & Average \\
\bottomrule
\end{tabular}
\end{center}

\subsubsection*{External Inquiries (EQ)}
\begin{center}
\begin{tabular}{clc}
\toprule
\textbf{\#} & \textbf{Function} & \textbf{Complexity} \\
\midrule
1 & View Attendance History & Low \\
2 & View Assigned Tasks (Article) & Low \\
3 & View All Articles List (Admin) & Low \\
4 & View Leave Request Status & Low \\
5 & View Task Progress (Admin) & Average \\
6 & View Notification History & Low \\
\bottomrule
\end{tabular}
\end{center}

\subsubsection*{Internal Logical Files (ILF)}
\begin{center}
\begin{tabular}{clc}
\toprule
\textbf{\#} & \textbf{Data Store (MySQL Table)} & \textbf{Complexity} \\
\midrule
1 & admins & Low \\
2 & articles (with face\_descriptor) & High \\
3 & attendance & High \\
4 & tasks & Average \\
5 & task\_assignments & Average \\
6 & progress\_log & Average \\
7 & leave\_requests & Average \\
8 & notifications & Low \\
9 & settings & Low \\
\bottomrule
\end{tabular}
\end{center}

\subsubsection*{External Interface Files (EIF)}
\begin{center}
\begin{tabular}{clc}
\toprule
\textbf{\#} & \textbf{External System} & \textbf{Complexity} \\
\midrule
1 & face-api.js (Face Recognition) & High \\
2 & Browser Geolocation API (GPS) & Average \\
3 & Nodemailer (SMTP Email) & Low \\
4 & Chart.js (Data Visualization) & Average \\
\bottomrule
\end{tabular}
\end{center}

\subsection{Step 2: Calculate Unadjusted Function Points (UFP)}

The standard weighting factors are:

\begin{center}
\begin{tabular}{lccc}
\toprule
\textbf{Type} & \textbf{Low} & \textbf{Average} & \textbf{High} \\
\midrule
External Inputs (EI) & 3 & 4 & 6 \\
External Outputs (EO) & 4 & 5 & 7 \\
External Inquiries (EQ) & 3 & 4 & 6 \\
Internal Logical Files (ILF) & 7 & 10 & 15 \\
External Interface Files (EIF) & 5 & 7 & 10 \\
\bottomrule
\end{tabular}
\end{center}

\bigskip
\noindent\textbf{Calculation:}

\begin{center}
\begin{tabular}{l c c c c}
\toprule
\textbf{Type} & \textbf{Low} & \textbf{Average} & \textbf{High} & \textbf{Total} \\
\midrule
EI & $4 \times 3 = 12$ & $5 \times 4 = 20$ & $1 \times 6 = 6$ & \textbf{38} \\
EO & $1 \times 4 = 4$ & $2 \times 5 = 10$ & $3 \times 7 = 21$ & \textbf{35} \\
EQ & $5 \times 3 = 15$ & $1 \times 4 = 4$ & --- & \textbf{19} \\
ILF & $3 \times 7 = 21$ & $4 \times 10 = 40$ & $2 \times 15 = 30$ & \textbf{91} \\
EIF & $1 \times 5 = 5$ & $2 \times 7 = 14$ & $1 \times 10 = 10$ & \textbf{29} \\
\midrule
\multicolumn{4}{r}{\textbf{Unadjusted Function Points (UFP)}} & $\mathbf{212}$ \\
\bottomrule
\end{tabular}
\end{center}

\subsection{Step 3: Value Adjustment Factor (VAF)}

The VAF is calculated using 14 General System Characteristics (GSCs), each rated 0--5.

\begin{center}
\begin{tabular}{clc}
\toprule
\textbf{\#} & \textbf{Characteristic} & \textbf{Rating (0--5)} \\
\midrule
1 & Data Communications & 4 \\
2 & Distributed Data Processing & 2 \\
3 & Performance & 4 \\
4 & Heavily Used Configuration & 2 \\
5 & Transaction Rate & 3 \\
6 & Online Data Entry & 5 \\
7 & End-User Efficiency & 4 \\
8 & Online Update & 4 \\
9 & Complex Processing & 5 \\
10 & Reusability & 3 \\
11 & Installation Ease & 3 \\
12 & Operational Ease & 4 \\
13 & Multiple Sites & 2 \\
14 & Facilitate Change & 3 \\
\midrule
 & \textbf{Total Degree of Influence (TDI)} & \textbf{48} \\
\bottomrule
\end{tabular}
\end{center}

\[
\text{VAF} = 0.65 + (0.01 \times \text{TDI}) = 0.65 + (0.01 \times 48) = \mathbf{1.13}
\]

\subsection{Step 4: Adjusted Function Points (AFP)}

\[
\text{AFP} = \text{UFP} \times \text{VAF} = 212 \times 1.13 = \mathbf{239.56 \approx 240~FP}
\]

\subsection{Step 5: Effort and Cost Estimation}

Using industry averages for a web application project:
\begin{itemize}[noitemsep]
    \item Productivity rate: \textbf{10 FP per person-month} (moderate complexity MERN-equivalent stack)
    \item Average developer cost: \textbf{\textrm{INR} 40,000 per month} (academic/intern-level team)
\end{itemize}

\[
\text{Effort} = \frac{\text{AFP}}{\text{Productivity}} = \frac{240}{10} = \mathbf{24.0~\text{person-months}}
\]

With a team of 4 developers:
\[
\text{Duration} = \frac{24.0}{4} = \mathbf{6.0~\text{months}}
\]

\[
\text{Total Cost} = 24.0 \times 40{,}000 = \mathbf{\textrm{INR}~9,60,000}
\]

% ═══════════════════════════════════════════════════════════
% PART B: USE CASE POINTS
% ═══════════════════════════════════════════════════════════
\section{Use Case Points (UCP) Estimation}

\subsection{Overview}
The Use Case Points method, proposed by Gustav Karner in 1993, estimates effort based on the complexity of actors and use cases, adjusted by technical and environmental factors.

\subsection{Step 1: Unadjusted Actor Weight (UAW)}

Actors are classified by their interaction type:
\begin{center}
\begin{tabular}{llcc}
\toprule
\textbf{Actor} & \textbf{Type} & \textbf{Weight} & \textbf{Count} \\
\midrule
Admin (Chartered Accountant) & Complex (GUI-based) & 3 & 1 \\
Article Trainee & Complex (GUI-based) & 3 & 1 \\
face-api.js (Face Recognition) & Average (library/protocol) & 2 & 1 \\
Browser Geolocation API (GPS) & Simple (API) & 1 & 1 \\
Nodemailer (Email SMTP) & Simple (API) & 1 & 1 \\
Chart.js (Visualization) & Simple (API) & 1 & 1 \\
\midrule
\multicolumn{3}{r}{\textbf{UAW}} & $\mathbf{11}$ \\
\bottomrule
\end{tabular}
\end{center}

\subsection{Step 2: Unadjusted Use Case Weight (UUCW)}

Use cases are classified by transaction count:
\begin{center}
\begin{tabular}{llcc}
\toprule
\textbf{Use Case} & \textbf{Complexity} & \textbf{Weight} & \textbf{Trans.} \\
\midrule
Login / Logout & Simple ($\leq$3) & 5 & 3 \\
Add Article Trainee & Average (4--7) & 10 & 5 \\
Edit/Remove Article & Simple & 5 & 3 \\
Mark Attendance (Face + GPS) & Complex ($>$7) & 15 & 12 \\
View Attendance History & Simple & 5 & 2 \\
Create \& Assign Task & Average & 10 & 6 \\
Update Task Progress & Average & 10 & 5 \\
Apply for Leave & Simple & 5 & 3 \\
Approve/Reject Leave & Simple & 5 & 3 \\
Generate Attendance Report & Complex & 15 & 9 \\
Generate Performance Report & Complex & 15 & 8 \\
View Analytics Dashboard & Average & 10 & 7 \\
Configure System Settings & Simple & 5 & 2 \\
\midrule
\multicolumn{2}{r}{\textbf{UUCW}} & $\mathbf{115}$ & \\
\bottomrule
\end{tabular}
\end{center}

\subsection{Step 3: Unadjusted Use Case Points (UUCP)}

\[
\text{UUCP} = \text{UAW} + \text{UUCW} = 11 + 115 = \mathbf{126}
\]

\subsection{Step 4: Technical Complexity Factor (TCF)}

\begin{center}
\begin{tabular}{clcc}
\toprule
\textbf{T\#} & \textbf{Factor} & \textbf{Wt.} & \textbf{Rating (0--5)} \\
\midrule
T1 & Distributed System & 2 & 3 \\
T2 & Response Time / Performance & 1 & 4 \\
T3 & End-User Efficiency & 1 & 4 \\
T4 & Complex Internal Processing & 1 & 5 \\
T5 & Reusability & 1 & 3 \\
T6 & Easy to Install & 0.5 & 3 \\
T7 & Easy to Use & 0.5 & 4 \\
T8 & Portability & 2 & 3 \\
T9 & Easy to Change & 1 & 3 \\
T10 & Concurrency & 1 & 3 \\
T11 & Security Features & 1 & 5 \\
T12 & Third-Party Access & 1 & 4 \\
T13 & Training Required & 1 & 2 \\
\midrule
\multicolumn{3}{r}{\textbf{TFactor} $= \sum(\text{Wt.} \times \text{Rating})$} & \textbf{48.5} \\
\bottomrule
\end{tabular}
\end{center}

\[
\text{TCF} = 0.6 + (0.01 \times 48.5) = \mathbf{1.085}
\]

\subsection{Step 5: Environmental Complexity Factor (ECF)}

\begin{center}
\begin{tabular}{clcc}
\toprule
\textbf{E\#} & \textbf{Factor} & \textbf{Wt.} & \textbf{Rating (0--5)} \\
\midrule
E1 & Familiarity with Project & 1.5 & 3 \\
E2 & Application Experience & 0.5 & 2 \\
E3 & OOP Experience & 1 & 4 \\
E4 & Lead Analyst Capability & 0.5 & 3 \\
E5 & Motivation & 1 & 4 \\
E6 & Stable Requirements & 2 & 4 \\
E7 & Part-Time Workers & $-1$ & 3 \\
E8 & Difficult Programming Language & $-1$ & 1 \\
\midrule
\multicolumn{3}{r}{\textbf{EFactor} $= \sum(\text{Wt.} \times \text{Rating})$} & \textbf{19.0} \\
\bottomrule
\end{tabular}
\end{center}

\[
\text{ECF} = 1.4 + (-0.03 \times 19) = 1.4 - 0.57 = \mathbf{0.83}
\]

\subsection{Step 6: Adjusted Use Case Points (UCP)}

\[
\text{UCP} = \text{UUCP} \times \text{TCF} \times \text{ECF} = 126 \times 1.085 \times 0.83 = \mathbf{113.4 \approx 113~UCP}
\]

\subsection{Step 7: Effort and Cost Estimation}

Using the standard productivity factor of \textbf{20 person-hours per UCP}:

\[
\text{Effort (hours)} = 113.4 \times 20 = \mathbf{2{,}268~\text{person-hours}}
\]

\noindent Converting to person-months (assuming 160 working hours/month):
\[
\text{Effort (months)} = \frac{2268}{160} = \mathbf{14.175 \approx 14.2~\text{person-months}}
\]

With a team of 4 developers:
\[
\text{Duration} = \frac{14.2}{4} \approx \mathbf{3.6~\text{months}}
\]

\[
\text{Total Cost} = 14.2 \times 40{,}000 = \mathbf{\textrm{INR}~5,68,000}
\]

% ═══════════════════════════════════════════════════════════
\section{Comparison of Results}
% ═══════════════════════════════════════════════════════════

\begin{center}
\begin{tabular}{lcc}
\toprule
\textbf{Parameter} & \textbf{FPA} & \textbf{UCP} \\
\midrule
Software Size & 240 FP & 113 UCP \\
Effort & 24.0 person-months & 14.2 person-months \\
Duration (4 developers) & 6.0 months & 3.6 months \\
Estimated Cost & INR 9,60,000 & INR 5,68,000 \\
\bottomrule
\end{tabular}
\end{center}

\medskip
\noindent The FPA method produces a higher estimate because it captures the complexity of all 9 MySQL database tables and the heavy internal processing involved in biometric face recognition (face descriptor comparison using Euclidean distance) and GPS-based geofencing (Haversine distance calculation). These internal computations significantly inflate the ILF and EI counts.

\noindent The UCP method yields a more conservative figure as it focuses on user-facing use cases and adjusts for the team's environment. The high ECF adjustment (due to stable requirements from a single-client CA firm) reduces the final estimate.

\noindent A realistic estimate for this project lies between the two bounds, approximately \textbf{19 person-months} of total effort, translating to roughly \textbf{4.7 months} of development with a 4-member team.

% ═══════════════════════════════════════════════════════════
\section{Conclusion}
% ═══════════════════════════════════════════════════════════
This assignment demonstrated two complementary approaches to software cost estimation for the CA Firm Attendance and Article Management System. Function Point Analysis provided a comprehensive size measurement accounting for all data interactions, biometric processing, and system boundaries, yielding 240 adjusted function points. The Use Case Points method offered an effort-centric estimate grounded in use case complexity and team environment factors, producing 109 adjusted use case points.

Together, these techniques provide upper and lower bounds for project planning, enabling the development team to set realistic milestones, allocate resources proportionally, and manage stakeholder expectations with data-driven estimates.

\end{document}
